\subsection{Метод переразмещения целых подграфов}
Конфигурирование выбирает среду до исполнения; переразмещение изменяет назначения при изменении её состояния. Здесь определена архитектура будущей интеграции. Её экспериментальная реализация, обучение и ускорение не заявляются.

Связная задача имеет фиксированное разбиение на целые подграфы. Перенос сохраняет состав каждой части, поэтому отдельная вершина не назначается независимо. Переразбиение образует новые части и остаётся отдельной операцией.

Алгоритм М2.1 проверяет актуальность, область применимости и сопоставимость внешних команд, затем фиксирует их на один поиск. М2.2 формирует полные назначения и выбирает допустимый план. М2.3 сохраняет или пересматривает решение при обновлении данных, регистрируя основания выбора.

Пусть $a$ обозначает полный план, а $R(a)$ --- команды назначенных устройств, упорядоченные по убыванию. Через $T(a)$ обозначим расчётное время, через $N(a)$ --- число переносов относительно прежнего решения. Защищённому эталону соответствует критерий
\begin{equation}\label{eq:placement}
 \min_{\mathrm{lex}}\bigl(R(a),T(a),N(a),\operatorname{id}(a)\bigr).
\end{equation}
Первое сравнение в~\eqref{eq:placement} является лексикографическим минимаксным. Идентификатор устраняет оставшиеся равенства. Команды порядковые: меньшая предпочтительнее, но не означает меньшее число секунд или вероятность. Без применимого прогноза действует порядок $(T,N,\operatorname{id})$.

Каждый план проходит проверку ресурсов, информационных зависимостей, передач и срока всей связной задачи. Прогноз не ослабляет ограничений. Результат --- выбранный полный план с протоколом, не рекомендация другому методу. При ограниченном поиске вывод относится только к проверенным планам; неуспех поиска не доказывает невозможности размещения.

\subsection{Модуль прогнозирования и калибровка}
Предусматриваются три оценивателя. Первый оценивает характеристики конфигурации, второй --- риск переполнения Fog, третий --- состояние источника данных. Возможные признаки включают p95, потери, $H(200)$, состояние датчика, интенсивность и занятость ресурсов.

Каждый оцениватель требует определённых целевой переменной, горизонта и области применения. Для отказов необходимы телеметрия и размеченные события; ручное отключение датчика не заменяет такие данные. Для переполнения должны учитываться очередь и обслуживание, а не только интенсивность.

Выход интерфейса выбора --- сопоставимые порядковые команды. Физические оценки и риски преобразуются в команды только по отдельно заданному и откалиброванному правилу. Потребуются отложенные конфигурации, проверка порядка предпочтений и критерий отказа от неприменимого прогноза. Ни одна из этих калибровок здесь не выполнена.

\begin{figure}[htbp]\centering
\begin{tikzpicture}[box/.style={draw,rounded corners,align=center,text width=88mm,minimum height=6mm,font=\small},>=Stealth]
\node[box] (a) {Снимок среды / Infrastructure snapshot};
\node[box,below=2mm of a] (b) {Будущие оцениватели / Future estimators};
\node[box,below=2mm of b] (c) {Применимость команд / Command applicability};
\node[box,below=2mm of c] (d) {Поиск целых подграфов / Whole-subgraph search};
\node[box,below=2mm of d] (e) {Проверка полного плана / Complete-plan verification};
\node[box,below=2mm of e] (f) {Выбор по $(R,T,N,\operatorname{id})$ / Selection};
\draw[->] (a)--(b);\draw[->] (b)--(c);\draw[->] (c)--(d);\draw[->] (d)--(e);\draw[->] (e)--(f);
\end{tikzpicture}
\caption{Архитектура будущей интеграции; прогнозный контроллер не выполнялся.\engcap{Architecture of the proposed integration; the forecast-driven controller has not been executed.}}\label{fig:placement}
\end{figure}

\subsection{Связь с моделью вычислительной среды}
Модель должна предоставлять данные для будущего обучения и служить стендом проверки метода. Каждая запись связывает снимок, входной поток, конфигурацию и наблюдаемые исходы. Обучающая и проверочная части разделяются до настройки оценивателя.

Схема рисунка~\ref{fig:placement} отделяет прогноз от допустимости. На одинаковых входах потребуется сопоставить базовый режим, контролируемые порядковые команды и обученный оцениватель. Синтетические команды проверяют алгоритм выбора, но не полезность действующей нейросети.

Перенос на одинаковый узел не гарантирует выигрыша. Результат зависит от неоднородности, очередей, каналов и стоимости переноса. Горизонт упреждения определяет окно подготовки, не гарантированный резерв. Разгрузка Fog в Cloud и оценка ускорения остаются дальнейшей работой.
